\cvsection{Expériences Académiques}

\cvexperience{Été 2026 \\ (4 mois)}{Chargé de cours \newline GEN430 -- Circuits logiques}
{Université de Sherbrooke}{Sherbrooke}
{Enseignant au baccalauréat en génie électrique}
{\vspace{-12pt}
\begin{itemize}
	\item Cours magistraux / Séminaires de consolidation des concepts
	\item Conception de machines à états de Moore et de Mealy
	\item Programmation de FPGA en VHDL et Verilog
\end{itemize}}

\cvexperience{Été 2026 \\ (4 mois)}{Chargé de cours \newline GEN420 -- Mathématiques des circuits logiques}
{Université de Sherbrooke}{Sherbrooke}
{Enseignant au baccalauréat en génie électrique}
{\vspace{-12pt}
\begin{itemize}
	\item Cours magistraux / Séminaires de consolidation des concepts
	\item Introduction à l'Algèbre booléen 
	\item Utilisation des tables de Karnaugh
	\item Conception de circuits logiques
\end{itemize}}

\cvexperience{Été 2025 \\ (4 mois)}{Auxiliaire d'enseignement \newline GEI816 -- Vérification fonctionnelle de systèmes numériques mixtes}
{Université de Sherbrooke}{Sherbrooke}
{Support technique pour les étudiants du baccalauréat en génie informatique}
{\vspace{-12pt}
\begin{itemize}
	\item Techniques de vérification des circuits intégrés
	\item Correction et rétroaction sur le travail (rapports, examens)
	\item Consultation pour la préparation aux évaluations
\end{itemize}}

\cvexperience{Automne 2025 \\ Été 2025 \\ (8 mois)}{Auxiliaire d'enseignement \newline GEL402 -- Conception d'un système numérique}
{Université de Sherbrooke}{Sherbrooke}
{Support technique pour les étudiants du baccalauréat en génie électrique}
{\vspace{-12pt}
\begin{itemize}
	\item Programmation de microprocesseurs et de FPGA
	\item Protocoles de communication (UART, SPI)
	\item Développement et mise à jour du matériel de cours
	\item Support en laboratoire (travaux pratiques) pour les étudiants
	% \item Tutorat extra-curriculaire non-payé pour les étudiants en difficulté
	\item Développement de nouveaux barèmes de corrections
\end{itemize}}
